% BevoTag System Design Report, memo-style report with decimal section numbers.
% Formatting follows resources/ECE Style Guide.pdf and resources/Template for Memo Style Reports.doc;
% content follows resources/System Design Report description.pdf and rubric (paths relative to this file).
% Preamble matches docs/prior-art/main.tex except for the decimal heading numbers this assignment requires.
% Build: pdflatex main.tex (twice).
\documentclass[12pt,letterpaper]{article}

\usepackage[T1]{fontenc}
\usepackage{newtxtext,newtxmath}              % 12-pt Times
\usepackage[margin=1in]{geometry}             % 1-in margins on all sides
\usepackage{setspace}
\setstretch{1.43}                             % Word "1.5 lines" for 12-pt Times (20.7-pt baselines)
\usepackage[document]{ragged2e}               % do not justify the right margin
\usepackage{graphicx}
\usepackage{xcolor}
\usepackage{array,tabularx,booktabs}
\usepackage{enumitem}
\usepackage{caption}
\usepackage{float}
\usepackage{flafter}                          % a table never appears before the text that cites it
\usepackage{placeins}
\usepackage{etoolbox}
\usepackage[nobottomtitles*]{titlesec}        % never leave a heading on the last line of a page
\usepackage{varioref}
\usepackage[hidelinks]{hyperref}
\usepackage{url}
\renewcommand{\UrlFont}{\normalfont}
\def\UrlBreaks{\do\/\do\-\do\_\do\=\do\?}  % no break after "." (avoids "doi: 10." line ends)

% Block paragraphs: no indentation, one empty line between paragraphs.
\setlength{\parindent}{0pt}
\setlength{\parskip}{\baselineskip}
\renewcommand{\bottomtitlespace}{0.05\textheight}

% Ragged-right text without hyphenated line breaks.
\hyphenpenalty=10000
\exhyphenpenalty=10000

% Headings (Style Guide 2.3) with the decimal numbering this assignment requires:
%   1st level: "1.0 INTRODUCTION" ALL CAPS, bold   2nd level: "3.1 Design Concept" bold
%   3rd level: "3.1.1 Tag Hardware" bold italic. References and appendixes use starred (unnumbered) forms.
\setcounter{secnumdepth}{3}
\renewcommand{\thesection}{\arabic{section}.0}
\renewcommand{\thesubsection}{\arabic{section}.\arabic{subsection}}
\renewcommand{\thesubsubsection}{\thesubsection.\arabic{subsubsection}}
\titleformat{\section}{\normalfont\bfseries}{\thesection}{0.5em}{\MakeUppercase}
\titleformat{\subsection}{\normalfont\bfseries}{\thesubsection}{0.5em}{}
\titleformat{\subsubsection}{\normalfont\bfseries\itshape}{\thesubsubsection}{0.5em}{}
\titlespacing*{\section}{0pt}{0pt}{-\parskip}
\titlespacing*{\subsection}{0pt}{0pt}{-\parskip}
\titlespacing*{\subsubsection}{0pt}{0pt}{-\parskip}


% Captions (Style Guide 2.4): bold, centered, "Table 1." above tables, "Figure 1." below figures.
\captionsetup{font={normalsize,bf},labelfont=bf,labelsep=period,justification=centering,singlelinecheck=false}
\captionsetup[table]{position=above,skip=6pt}
\captionsetup[figure]{position=below,skip=6pt}
\renewcommand{\topfraction}{0.9}
\renewcommand{\bottomfraction}{0.8}
\renewcommand{\textfraction}{0.07}
\renewcommand{\floatpagefraction}{0.9}
\setcounter{topnumber}{3}
\setcounter{totalnumber}{4}
% One empty line between a graphic and the text around it.
\setlength{\textfloatsep}{\baselineskip}
\setlength{\floatsep}{\baselineskip}
\setlength{\intextsep}{\baselineskip}
\makeatletter
\g@addto@macro\float@endH{\vskip-\parskip\vskip\z@}
\makeatother
\setlength{\emergencystretch}{3em}
% No single line of a paragraph stranded at the top or bottom of a page.
\widowpenalty=10000
\clubpenalty=10000

% Indented lists that continue the preceding paragraph (Style Guide 2.5).
\setlist[itemize]{leftmargin=0.5in,topsep=0pt,itemsep=0pt,parsep=0pt}
\setlist[enumerate]{leftmargin=0.5in,topsep=0pt,itemsep=0pt,parsep=0pt}

% Single-spaced 12-pt tables.
\newcolumntype{L}[1]{>{\RaggedRight\arraybackslash}p{#1}}
\newcolumntype{Y}{>{\RaggedRight\arraybackslash}X}
\newenvironment{spectable}{\setstretch{1.0}\renewcommand{\arraystretch}{1.1}}{}

% Page numbers centered at the bottom, on every page but the first.
\pagestyle{plain}

\pretocmd{\section}{\FloatBarrier}{}{}
% Pointer to a graphic that does not share its first reference's page (Style Guide 2.4).
\vrefwarning
\renewcommand{\reftextcurrent}{}
\renewcommand{\reftextafter}{on the next page}
\renewcommand{\reftextfaceafter}{on the next page}
\begin{document}
\thispagestyle{empty}

%----------------------------------------------------------------------------------------------------
% MEMO HEADING (Style Guide 2.2): must match the template exactly. DATE is the due date.
%----------------------------------------------------------------------------------------------------
\begingroup
\setstretch{1.0}
\newcommand{\fieldgap}{\\[6.2pt]}
\noindent\begin{tabular}{@{}p{1.1in}@{}p{\dimexpr\linewidth-1.1in\relax}@{}}
\textbf{DATE:} & December 4, 2026 \fieldgap
\textbf{TO:} & Dr.\ Haris Vikalo \\ & Mugdha Jadhao \fieldgap
\textbf{FROM:} & Ethan Chua \\ & Esteban Martinez \\ & Li Shi \\ & Andrew Yin \\ & Abhirit Das \fieldgap
\textbf{SUBJECT:} & System Design Report for BevoTag: TODO subtitle \\
\end{tabular}
\endgroup

%----------------------------------------------------------------------------------------------------
\section{Introduction}\label{sec:intro}
%----------------------------------------------------------------------------------------------------
% Purpose -> context (Problem Statement, Prior Art Report, Design Review) -> summary of contents.
% An engineer new to BevoTag should grasp the problem, the design approach, and the 464 plan from
% this section alone. No design details here (TA complaint on the Project Definition).
TODO.

%----------------------------------------------------------------------------------------------------
\section{Design Project Background}\label{sec:background}
%----------------------------------------------------------------------------------------------------
% Summary of the Problem Statement, updated; fix every TA comment on reused text.
% Must cover: the design problem, specifications (for a technical reader), and deliverables.
% Refer to Appendix A for standards.
TODO overview paragraph.

\subsection{Design Problem}
TODO.

\subsection{Specifications}
TODO (table of measurable specs; reference Appendix~A).

\subsection{Deliverables}
TODO.

%----------------------------------------------------------------------------------------------------
\section{System Design}\label{sec:design}
%----------------------------------------------------------------------------------------------------
% Updated and expanded Design Review material. Needs an overview paragraph before 3.1.
TODO overview paragraph.

\subsection{Design Concept}\label{sec:concept}
% Start with a system-level overview and the block diagram (Figure 1), explained in the text.
% For EVERY subsystem: how it works, inputs/outputs to other blocks, why chosen
% (alternatives rejected, test results, calculations/research, which spec it meets).
% Include the user interface as a subsystem.
TODO system overview and block diagram.

% \begin{figure}[htbp]
%   \centering
%   \includegraphics[width=\linewidth]{figures/block-diagram.png}
%   \caption{BevoTag System Block Diagram}
%   \label{fig:block}
% \end{figure}

\subsubsection{Tag Hardware}
TODO.

\subsubsection{Tag Firmware}
TODO.

\subsubsection{Survey and Fingerprint Database}
TODO.

\subsubsection{Location Server and Sensor Fusion}
TODO.

\subsubsection{User Interface}
TODO (2D/3D map application; meet the Problem Statement's UI requirements).

\subsection{Risk Reduction}\label{sec:risk}
% Comprehensive: methods, results, and analysis of every lab experiment this semester.
% Source material: docs/risk-reduction-plan.md, docs/individual-risk-reduction-demos.md,
% and each member's answers in answers/*.md (Q1--Q4).
% Owner names stay in these comments only; the report text is third person.
TODO overview paragraph (which risks were tested and why).

\subsubsection{Wi-Fi Data Collection and Antenna Comparison}
% Team experiments from docs/risk-reduction-plan.md: room hit rate, median error,
% PCB vs. omnidirectional antenna RSSI variation.
TODO.

\subsubsection{Router Localization Firmware}
% Owner: Ethan Chua. Firmware that locates a router by logging the strongest signal strength.
TODO.

\subsubsection{Access Point Inventory for EER and ECJ}
% Owner: Abhirit Das. List of the WAPs in EER and ECJ.
TODO.

\subsubsection{Preliminary Access Point Locations}
% Owner: Andrew Yin. Preliminary list of WAP locations.
TODO.

\subsubsection{Survey Data Parsing and Display Software}
% Owner: Esteban Martinez. Display software for parsing data collections.
TODO.

\subsubsection{3D Building Model and Signal Map}
% Owner: Li Shi. Web app 3D model and signal map.
TODO.

\subsection{Test Plan}\label{sec:testplan}
% Key test(s) per block of Figure 1, plus system-level test(s).
% For each: procedure, data collected, success criterion (tie to the specs in 2.2).
TODO.

%----------------------------------------------------------------------------------------------------
\section{Project Management}\label{sec:pm}
%----------------------------------------------------------------------------------------------------
% Overview paragraph required. Each subsection summarizes one appendix (B--E) and refers to it.
TODO overview paragraph.

\subsection{Project Activities}
TODO (summarize the WBS in Appendix~B, organized around major tasks).

\subsection{Task Assignments}
TODO (summarize the LRC in Appendix~C; rationale: equitable load and qualifications).

\subsection{Project Schedule}
TODO (summarize the Gantt chart in Appendix~D chronologically, with milestone dates).

\subsection{Project Budget}
TODO (summarize the BOM in Appendix~E; material cost goals, organized to match the table).

%----------------------------------------------------------------------------------------------------
\section{Conclusion}\label{sec:conclusion}
%----------------------------------------------------------------------------------------------------
% For a manager who reads nothing else: purpose, main points, and a clear assessment of the
% team's readiness for EE 464 and its prospects for success. Keep it short.
TODO.

%----------------------------------------------------------------------------------------------------
% REFERENCES (IEEE style, required)
%----------------------------------------------------------------------------------------------------
\clearpage
\section*{REFERENCES}
\begingroup
\renewcommand{\section}[2]{}%   suppress thebibliography's own heading
\begin{thebibliography}{10}
\bibitem{todo}
TODO: reuse entries from docs/prior-art/main.tex and docs/project-definition/main.tex.
\end{thebibliography}
\endgroup

%----------------------------------------------------------------------------------------------------
% APPENDIXES A--E (all required). Divider page with the title one-third down, content on the next
% page; page numbers continue from the body (Style Guide 2.7).
%----------------------------------------------------------------------------------------------------
\newcommand{\appendixdivider}[1]{%
  \clearpage\vspace*{0.20\textheight}\begin{center}\textbf{#1}\end{center}\clearpage}

\appendixdivider{APPENDIX A -- RELEVANT STANDARDS}
% Updated from the Project Definition's appendix; address every TA comment on it.
TODO.

\appendixdivider{APPENDIX B -- WORK BREAKDOWN STRUCTURE}
% All activities and major tasks from submission of this report to the EE 464 prototype.
TODO.

\appendixdivider{APPENDIX C -- LINEAR RESPONSIBILITY CHART}
% Responsibilities per team member for every WBS task.
TODO.

\appendixdivider{APPENDIX D -- GANTT CHART}
% Every WBS activity and milestone, from right after submission to the end of EE 464.
TODO.

\appendixdivider{APPENDIX E -- BILL OF MATERIALS}
% Major items only: description, quantity, performance criteria (part numbers if known).
TODO.

\end{document}
